\documentclass[10pt,a4paper]{amsart}
\usepackage[margin=19mm]{geometry}
\usepackage{amsmath,amssymb,mathtools}
\usepackage[scr=rsfs,cal=euler]{mathalfa}
\usepackage{xfrac}
\usepackage[unicode,hidelinks,bookmarks=false]{hyperref}
\newcommand{\ie}{i.e.\ }
\newcommand{\cf}{cf.\ }
\newcommand{\resp}{respectively\ }
\newcommand{\Subsection}{\subsection}
\providecommand{\nobreakdash}{\nobreak\hbox{-}}
\setlength{\emergencystretch}{2em}
\multlinegap0pt
\usepackage{mdframed}
\usepackage{mdframed}
\usepackage{mdframed}
\usepackage{mdframed}
\usepackage{amssymb}
\usepackage{csquotes}
\usepackage{enumerate}
\usepackage{float}
\usepackage{mathtools}
\usepackage{mdframed}
\usepackage{tikz}
\newtheorem{proposition}{Proposition}
\title{Two eggs any style \\ generalizing egg-drop experiments}
\author{Harold R. Parks, Dean C. Wills}
\date{Source: arXiv:2208.09576v1 (CC BY 4.0)}
\begin{document}
\maketitle

\noindent\emph{Source and licence.} Two eggs any style \\ generalizing egg-drop experiments. Version arXiv:2208.09576v1 (CC BY 4.0), distributed by arXiv under the Creative Commons Attribution 4.0 International licence, http://creativecommons.org/licenses/by/4.0/, as recorded in the arXiv licence metadata for this exact version and shown on the abstract page https://arxiv.org/abs/2208.09576v1. Full licence notice: \texttt{licenses/arxiv-2208.09576-CC-BY-4.0.txt}.\par\smallskip

\noindent\textit{Editorial scope.} Counted unit: the article's proposition normal-non (source0.tex lines 521--524). The article's complete original proof of it is delivered with it (source0.tex lines 525--615, one \texttt{proof} environment of 91 lines). No mathematics is written, repaired or re-derived: the statement and the proof are the article's own lines, in the article's own notation, and no retained line is substituted or re-flowed.  No further source-local context is needed: every cross-reference of every retained line resolves inside the two delivered spans.  Nothing was omitted from any retained span, so the package carries no omission record.  No retained line contains a citation, so the wrapper carries no bibliography and no bibliography entry.  No retained line contains a figure environment, an image-inclusion command or drawing code, so no image is delivered and the package declares no asset dependency.  Editorial formatting only, in addition: an AMS \texttt{amsart} preamble that restates the article's own declarations and macros used by the retained lines, namely the environments \texttt{proposition} on separate counters, together with the standard packages those lines need.  The retained environments print in this document as Proposition 1 in order of appearance, whereas the article prints the same items with the numbering of its own full document.  The complete native source of the article as distributed in the arXiv e-print for this exact version is bundled unchanged as \texttt{sources/129674/source0.tex}.
\begin{proposition}\label{normal-non} 
Any algorithm for solving any of the egg-drop  problem variations  can be modified 
so as to be non-redundant and to never require more experiments than the original algorithm. 
\end{proposition}

\begin{proof}
We will show that any algorithm that requires 
an experiment for which the outcome is guaranteed can be modified
so that that experiment is either omitted altogether or is
replaced by another experiment that does not have a 
guaranteed outcome. In any case, no more  experiments will be required
than were required by the original algorithm. Further modifications can be made 
until all experiments with guaranteed outcomes are eliminated.
This process is explained in detail below:

In the tree representation,
an experiment for which the outcome is guaranteed
will correspond to an inner node with only one child---recall
inaccessible nodes have all been removed. We consider such an inner 
node of least depth, say  $d$. After all nodes at depth $d$ have
been processed, we progress to nodes of greater depth.
When modifications are made to the algorithm, new inaccessible nodes
might be created, so when dealing with a particular node, removing
inaccessible descendants  is the first thing to do.

\medskip
For a node with a guaranteed outcome, let $x$ be the floor from which the
egg is to be dropped
and let $[y,z]$ be the solution range associated with that node. 
Since the situations are different, we will need to
consider separately the case when the guaranteed outcome is failure
and when the guaranteed outcome is success.

If $y=z$, then the solution is already known to be $y$. Since $y$ is known 
to be the solution, the node we are considering should instead be a leaf with $y$ inside 
the square to indicate that $y$ is the solution.
The entire left or  right child subtree below 
the node should be eliminated. The node is now a leaf, and the height of the
tree has not been increased.

Assume now that $y<z$, so  the solution $s$ could be any integer in $[y,z]$, 
and the algorithm cannot yet terminate.

\medskip
\noindent
\textbf{Guaranteed failure.}  Since dropping on $x$ will give no new information
and will lose an egg, this experiment could simply be omitted. Instead,
let us note that for failure to be guaranteed, it must hold that $x>z$.
By dropping the egg on $z$ instead of $x$, we might have a success and that
would tells us that the solution is $z$. If the egg breaks, then we are no
worse off than before. We have gained the information that the egg breaks on $z$,
but we can ignore the additional information for now and follow the
original algorithm.
For the tree, the label of the  node should be changed to $z$,
a right-child leaf labeled $z$ should be added, and the left-child 
subtree is unchanged. The node now has two children, and 
the height of the tree has not been increased. 

\medskip
\noindent
\textbf{Guaranteed success.}
Because success is guaranteed, it must hold that $x\leq y$.
The drop on $x$ is guaranteed to result in the egg surviving, 
so it may be  beneficial  in that  replacement or bonus eggs will be acquired. But since
no new information is gained, 
at least one additional experiment beyond the drop on $x$ will be required.  

Consider what happens if  the drop is made from $y+1$ instead of from $x$.
If the egg breaks, then the solution is known to be $y$. No additional 
experiments are required and the modified algorithm can terminate.
If the egg survives, then the benefit of acquiring replacement or bonus 
eggs is achieved as would have happened with a drop on $x$. 
We have gained the additional information that the egg survives when 
dropped on $y+1$, but we may ignore that information for now and 
simply follow the original algorithm.
For the tree, the label of the node should be changed to $y+1$,
a left-child leaf labeled $y$ should be added, and the right-child
subtree is unchanged. The node now has two children, and 
the height of the tree has not been increased. 

\medskip
The information that the egg broke when dropped on $z$
or  that the egg  survived when dropped on $y+1$ was 
temporarily ignored when we chose to leave the  subtree unchanged. Nonetheless
the information will affect the solution range
of descendant nodes.
We might well find newly inaccessible nodes and   new instances of guaranteed outcome
experiments in the   subtree.  Such issues occur at depth greater than $d$
and are to be dealt with after all depth $d$ nodes 
have been processed. When all nodes at depth $d$ have been processed, they are
all either leaves or inner nodes with two children. 
%\todo[inline]{Added ``then"}
Then when all nodes in the tree have
been processed the tree will be a full binary tree with height no greater than the 
original tree.
\end{proof}

\end{document}
